\section{Procesos y Subprocesos}

De acuerdo con la estructura de una organización enfocada en el mantenimiento industrial, un sistema de software efectivo se articula en torno a macroprocesos que van desde la gestión de la información fundamental hasta la ejecución de estrategias de mantenimiento \parencite{ibm:2024, stum:2018}. Herramientas como los sistemas EAM o CMMS son fundamentales para centralizar y automatizar estas funciones, con el fin de optimizar el rendimiento y la vida útil de los activos \parencite{fluke:2020, sap:sf, tractian:2025a}.

\subsection{Procesos de Soporte y Gestión de Información}

\begin{enumerate}
    \item \textbf{Gestión de Activos:} Proceso fundamental que consiste en crear y mantener un registro centralizado de todos los activos de la planta, incluyendo sus componentes, manuales técnicos, historial y garantías. Este repositorio único es la piedra angular de cualquier sistema de mantenimiento moderno \parencite{ftmaintenance:2019, manwinwin:2015, mazenko:2015, servicenow:sf}.
    \item \textbf{Gestión de Inventario (MRO):} Controla el stock de piezas de repuesto, herramientas y otros materiales (Maintenance, Repair, and Operations). Este módulo incluye la gestión de compras y la asignación de recursos a las órdenes de trabajo, siendo un componente central en la mayoría de las plataformas CMMS/EAM \parencite{ftmaintenance:2019, honeywell:2018, laurila:2017, stum:2018}.
    \item \textbf{Visualización e Interacción con Gemelo Digital:} Permite a los usuarios navegar por una representación 2D/3D de la planta y activar capas de información (eléctrica, mecánica, HSEQ). Al seleccionar un activo, se accede a su "hoja de vida" digital, que es una representación dinámica que emula el rendimiento de su contraparte física y puede incluir dashboards de salud en tiempo real, historial de intervenciones y vistas explosionadas de sus componentes \parencite{chen:2025, miao:2025, osullivan:2020}.
    \item \textbf{Generación de Informes y Análisis (KPIs):} El sistema procesa de forma continua los datos operativos y de mantenimiento para generar dashboards e informes automatizados. Estos informes presentan métricas clave para la gestión, como el Tiempo Medio Entre Fallos (MTBF), el Tiempo Medio de Reparación (MTTR) y la Efectividad General del Equipo (OEE), permitiendo un análisis basado en datos para la toma de decisiones gerenciales \parencite{fiix:2019, ftmaintenance:2019, hxgn:sf, obrien:sf}.
\end{enumerate}

\subsection{Procesos de Ejecución de Mantenimiento}

\begin{enumerate}
    \item \textbf{Ciclo de Mantenimiento Preventivo:} Este proceso proactivo se inicia cuando el Planificador de Mantenimiento programa órdenes de trabajo (OT) recurrentes, basándose en intervalos de tiempo fijos (ej., "inspección trimestral") o en métricas de uso del equipo. Una de las capacidades centrales de un CMMS es generar y asignar automáticamente estas órdenes de trabajo según el calendario o los umbrales de uso establecidos, asegurando que el mantenimiento se realice antes de que ocurra una falla \parencite{accruent:2025, collabit:2020, fmx:2019, honeywell:2018, laurila:2017, mazenko:2015}.
    \item \textbf{Ciclo de Mantenimiento Correctivo (Reactivo):} Este ciclo se activa tras la detección de un fallo y comienza con una \textbf{solicitud de trabajo (SR)}, que es la notificación inicial de una necesidad. El Supervisor revisa, aprueba y prioriza esta solicitud, convirtiéndola en una \textbf{orden de trabajo (OT) formal}. Posteriormente, el Planificador asigna los recursos necesarios, incluyendo el técnico y los repuestos. Al finalizar la reparación, el Técnico documenta el cierre, registrando información crucial para el historial del activo, como los códigos de falla, la descripción de la solución, las horas-hombre empleadas y los materiales utilizados \parencite{accruent:2025, fracttal:2024, landau:2023, stum:2018, tractian:2025b}.
    \item \textbf{Ciclo de Mantenimiento Predictivo (PdM):} Esta estrategia avanzada se fundamenta en la monitorización en tiempo real del estado de los activos mediante sensores IoT (ej. vibración, temperatura). Los Gemelos Digitales (DT), potenciados con modelos de Machine Learning (ML), procesan estos datos para pronosticar fallos inminentes y estimar la vida útil restante (RUL) de los componentes. El objetivo es transformar la gestión de mantenimiento de un enfoque reactivo a uno \textbf{proactivo}. El sistema genera una \textbf{alerta predictiva} que desencadena automáticamente una orden de trabajo, permitiendo una intervención justo a tiempo, maximizando la vida útil del activo y minimizando el tiempo de inactividad no planificado \parencite{chen:2025, ibm:2024, osullivan:2020, tractian:2025a}.
\end{enumerate}

\subsection{Procesos de Interoperabilidad}

\begin{enumerate}
    \item \textbf{Procesamiento Automatizado de Informes (LLM):} Para optimizar la gestión de datos no estructurados, el sistema utiliza un Large Language Model (LLM). Cuando un Contratista Externo sube un informe de trabajo (PDF, imagen), el LLM extrae la información clave (piezas, horas, descripción) y la presenta en un formato estandarizado para la aprobación del Supervisor. Esta capacidad se alinea con el uso de la Inteligencia Artificial Generativa (GenAI) para automatizar tareas complejas y repetitivas en el mantenimiento moderno \parencite{chen:2025, hxgn:sf}.
    \item \textbf{Sincronización con Sistemas Externos:} Se contempla la integración del sistema con plataformas empresariales existentes para asegurar la consistencia de los datos en toda la organización. La conexión con sistemas como ERP (Enterprise Resource Planning) es fundamental para unificar los flujos de información entre mantenimiento, finanzas y la cadena de suministro. Esta capacidad, denominada \textbf{"Interacción del Sistema"}, se realiza a través de APIs y es clave para un ecosistema EAM verdaderamente integrado \parencite{anylogic:sf, ibm:2024, laurila:2017, tractian:2025a, zhang:2025}.
\end{enumerate}
